\section{主定理}

\begin{theorem}[Alpöge の形式証明]\label{thm:grh}
\(q\ge 1\) を整数とし、\(\chi\) を法 \(q\) のディリクレ指標とする。
\begin{enumerate}
\item \(\RePart s>1/2\) ならば \(L(s,\chi)\ne0\) である。ただし \(\chi\) が主指標で
      \(s=1\) の場合を除く。
\item \(0<\RePart s<1\) において \(L(s,\chi)=0\) ならば、\(\RePart s=1/2\) である。
\item リーマンゼータ関数の零点は、自明零点 \(-2,-4,\ldots\) を除いて
      \(\RePart s=1/2\) 上にある。
\end{enumerate}
\end{theorem}

\noindent
ここで \(L(s,\chi)=\sum_{n\ge1}\chi(n)n^{-s}\) は
\(\RePart s>1\) からの有理型接続であり、\(\chi\) の原始性は仮定しない。
法 \(1\) の指標の \(L\) 関数は \(\zeta(s)\) である。
底本の五つの定理文のうち、ゼータ関数に関する二つは (3) を、
ディリクレ \(L\) 関数に関する三つは (1)、(2) および (2) の原始指標の場合を述べる。
\leansrc{Statements}{定理文}と
\leansrc{Main}{主定理の宣言}を参照されたい。

まず、素数冪和の片側上界から (1) を導く
（第\ref{sec:analytic}節）。次に、補助的な数体の無限素点に重さを割り当てることで、
この剰余をテータ微分の整除性を測る量として表す（第\ref{sec:data}節）。
第\ref{sec:lattices}節では、同じグラム行列式に対して
無限素点と有限素点の情報を合わせる方法を述べる。
第\ref{sec:upper-bound}節から第\ref{sec:determinant-estimates}節までで
比較に必要な上界、次元評価、捻りによる行列式比の評価をそれぞれ与え、
第\ref{sec:conclusion}節で定理を導く。

本稿は Alpöge の公表文 \cite{alpoege} に記された証明のうち、
Sernau らが整理した Lean 開発の可読版 \cite{grh-source} に拠る。
本文の L で始まる参照番号から、巻末にある底本の宣言を参照できる。
とくに局所テータ核の構成や補助捻りを得るための合同の評価には、底本の定理を引用する。
青字の「補足」は、本稿で読解のために付した注記である。
証明の着想と関連資料は、直後の概説にまとめた。
